\documentclass{article}
\usepackage[utf8]{inputenc}
\usepackage{hyperref}
\usepackage{listings}
\usepackage{xcolor}
\usepackage{enumitem}

\lstset{
    basicstyle=\ttfamily\small,
    breaklines=true,
    frame=single,
    backgroundcolor=\color{gray!5},
}

\title{Pretty-print SSH and PGP public keys with details}
\author{Marcos de Carvalho}
\date{2026-05-03}

\begin{document}

\section*{Pretty-print SSH and PGP public keys with details}
\textbf{Author:} Marcos de Carvalho \hfill \textbf{Date:} 2026-05-03

\noindent Reads SSH public key files and a GnuPG public keyring, then prints paths, SSH key text, PGP identities, long key IDs, fingerprints, and armored public key blocks in a formatted report.

\section*{Language}
bash

\section*{Category}
cryptography

\section*{Command}
\begin{lstlisting}
bash -c 'sshdir=${1:-$HOME/.ssh}; gpghome=${2:-$HOME/.gnupg}; printf "\033[1m=== SSH PUBLIC KEYS (%s) ===\033[0m\n" "$sshdir"; mapfile -d "" s < <(find "$sshdir" -type f \( -name "*.pub" -o -name authorized_keys \) -print0 2>/dev/null); ((${#s[@]})) || echo "No SSH public key files found."; for f in "${s[@]}"; do printf "\n\033[1;34m%s\033[0m\n" "$f"; ssh-keygen -lf "$f" 2>/dev/null | sed "s/^/  details: /"; sed "s/^/  key: /" "$f"; done; printf "\n\033[1m=== PGP PUBLIC KEYS (%s) ===\033[0m\n" "$gpghome"; if [ -e "$gpghome/pubring.kbx" ] || [ -e "$gpghome/pubring.gpg" ]; then GNUPGHOME="$gpghome" gpg --batch --list-public-keys --with-fingerprint --with-subkey-fingerprint --keyid-format long 2>/dev/null | sed "s/^/  /"; printf "\n\033[1m=== ARMORED PGP PUBLIC KEY BLOCKS ===\033[0m\n"; GNUPGHOME="$gpghome" gpg --batch --armor --export 2>/dev/null; else echo "No GnuPG public keyring found."; fi' _
\end{lstlisting}

\section*{Explanation}
The command uses a bash -c wrapper so both optional paths are supplied as trailing arguments. SSH\_DIR defaults to \textasciitilde{}/.ssh and is searched recursively for *.pub files and authorized\_keys files. Each matching SSH file is labeled, summarized with ssh-keygen -lf, and printed with indentation. GNUPG\_HOME defaults to \textasciitilde{}/.gnupg; when it contains pubring.kbx or pubring.gpg, GnuPG lists public keys with long key IDs, primary fingerprints, and subkey fingerprints, then exports all public keys as ASCII-armored blocks. The \_ placeholder acts as \$0 inside bash -c so \$1 and \$2 map cleanly to SSH\_DIR and GNUPG\_HOME.

\section*{Tags}
ssh, pgp, gpg, public-keys, fingerprints, cryptography, audit, pretty-print

\section*{Dependencies}
bash, find, ssh-keygen, sed, gpg


\section*{Arguments}
\begin{enumerate}
  \item \textbf{SSH\_DIR} (Optional): Directory to search recursively for SSH public key files. Defaults to \textasciitilde{}/.ssh. \\ Default: \textasciitilde{}/.ssh
  \item \textbf{GNUPG\_HOME} (Optional): GnuPG home directory whose public keyring will be listed and exported. Defaults to \textasciitilde{}/.gnupg. \\ Default: \textasciitilde{}/.gnupg
\end{enumerate}

\section*{Examples}
\begin{enumerate}
  \item \texttt{bash -c 'sshdir=\$\{1:-\$HOME/.ssh\}; gpghome=\$\{2:-\$HOME/.gnupg\}; printf "\textbackslash\{\}033[1m=== SSH PUBLIC KEYS (\%s) ===\textbackslash\{\}033[0m\textbackslash\{\}n" "\$sshdir"; mapfile -d "" s < <(find "\$sshdir" -type f \textbackslash\{\}( -name "*.pub" -o -name authorized\_keys \textbackslash\{\}) -print0 2>/dev/null); ((\$\{\#s[@]\})) || echo "No SSH public key files found."; for f in "\$\{s[@]\}"; do printf "\textbackslash\{\}n\textbackslash\{\}033[1;34m\%s\textbackslash\{\}033[0m\textbackslash\{\}n" "\$f"; ssh-keygen -lf "\$f" 2>/dev/null | sed "s/\textasciicircum{}/  details: /"; sed "s/\textasciicircum{}/  key: /" "\$f"; done; printf "\textbackslash\{\}n\textbackslash\{\}033[1m=== PGP PUBLIC KEYS (\%s) ===\textbackslash\{\}033[0m\textbackslash\{\}n" "\$gpghome"; if [ -e "\$gpghome/pubring.kbx" ] || [ -e "\$gpghome/pubring.gpg" ]; then GNUPGHOME="\$gpghome" gpg --batch --list-public-keys --with-fingerprint --with-subkey-fingerprint --keyid-format long 2>/dev/null | sed "s/\textasciicircum{}/  /"; printf "\textbackslash\{\}n\textbackslash\{\}033[1m=== ARMORED PGP PUBLIC KEY BLOCKS ===\textbackslash\{\}033[0m\textbackslash\{\}n"; GNUPGHOME="\$gpghome" gpg --batch --armor --export 2>/dev/null; else echo "No GnuPG public keyring found."; fi' \_} - Print SSH public keys from \textasciitilde{}/.ssh and PGP public keys from \textasciitilde{}/.gnupg
  \item \texttt{bash -c 'sshdir=\$\{1:-\$HOME/.ssh\}; gpghome=\$\{2:-\$HOME/.gnupg\}; printf "\textbackslash\{\}033[1m=== SSH PUBLIC KEYS (\%s) ===\textbackslash\{\}033[0m\textbackslash\{\}n" "\$sshdir"; mapfile -d "" s < <(find "\$sshdir" -type f \textbackslash\{\}( -name "*.pub" -o -name authorized\_keys \textbackslash\{\}) -print0 2>/dev/null); ((\$\{\#s[@]\})) || echo "No SSH public key files found."; for f in "\$\{s[@]\}"; do printf "\textbackslash\{\}n\textbackslash\{\}033[1;34m\%s\textbackslash\{\}033[0m\textbackslash\{\}n" "\$f"; ssh-keygen -lf "\$f" 2>/dev/null | sed "s/\textasciicircum{}/  details: /"; sed "s/\textasciicircum{}/  key: /" "\$f"; done; printf "\textbackslash\{\}n\textbackslash\{\}033[1m=== PGP PUBLIC KEYS (\%s) ===\textbackslash\{\}033[0m\textbackslash\{\}n" "\$gpghome"; if [ -e "\$gpghome/pubring.kbx" ] || [ -e "\$gpghome/pubring.gpg" ]; then GNUPGHOME="\$gpghome" gpg --batch --list-public-keys --with-fingerprint --with-subkey-fingerprint --keyid-format long 2>/dev/null | sed "s/\textasciicircum{}/  /"; printf "\textbackslash\{\}n\textbackslash\{\}033[1m=== ARMORED PGP PUBLIC KEY BLOCKS ===\textbackslash\{\}033[0m\textbackslash\{\}n"; GNUPGHOME="\$gpghome" gpg --batch --armor --export 2>/dev/null; else echo "No GnuPG public keyring found."; fi' \_ /etc/ssh \textasciitilde{}/.gnupg} - Print system SSH host public keys from /etc/ssh and the current user's GnuPG public keyring
\end{enumerate}

\section*{Output}
\begin{lstlisting}
\u001b[1m=== SSH PUBLIC KEYS (/home/user/.ssh) ===\u001b[0m

\u001b[1;34m/home/user/.ssh/id_ed25519.pub\u001b[0m
  details: 256 SHA256:examplefingerprint user@example (ED25519)
  key: ssh-ed25519 AAAAC3NzaC1lZDI1NTE5AAAAIExampleKey user@example

\u001b[1m=== PGP PUBLIC KEYS (/home/user/.gnupg) ===\u001b[0m
  pub   ed25519/0123456789ABCDEF 2026-05-03 [SC]
        0123 4567 89AB CDEF 0123  4567 0123 4567 89AB CDEF
  uid                 [ultimate] User Example <user@example>
  sub   cv25519/FEDCBA9876543210 2026-05-03 [E]
        FEDC BA98 7654 3210 FEDC  BA98 7654 3210 FEDC BA98

\u001b[1m=== ARMORED PGP PUBLIC KEY BLOCKS ===\u001b[0m
-----BEGIN PGP PUBLIC KEY BLOCK-----
...
-----END PGP PUBLIC KEY BLOCK-----
\end{lstlisting}

\section*{Notes}
\begin{itemize}
  \item The SSH section includes *.pub files and authorized\_keys files under SSH\_DIR.
  \item ssh-keygen -lf prints one fingerprint line per public key when a file contains multiple keys.
  \item The PGP section uses the selected GnuPG home directory and requires pubring.kbx or pubring.gpg to exist.
  \item Use a pager such as less -R when the exported PGP public key blocks are long.
  \item This command is intended for Bash because it uses mapfile and process substitution.
\end{itemize}

\section*{Warnings}
\begin{itemize}
  \item Public keys are not secret, but they often include names, email addresses, hostnames, or comments that may be sensitive in shared logs or screenshots.
  \item The command does not print private SSH or PGP key material.
  \item GnuPG may create or refresh local keyring metadata such as trustdb.gpg when inspecting a GnuPG home that lacks it.
\end{itemize}

\section*{See Also}
\begin{itemize}
  \item bulk-pgp-detach-sign-files
\end{itemize}

\section*{Status}
reviewed

\section*{Safety}
caution

\section*{Shell}
bash

\section*{Platforms}
linux, gnu-linux

\section*{Created At}
2026-05-03

\section*{Updated At}
2026-05-03

\end{document}
